\chapter{QUANTUM FRAMEWORKS}

This block focuses on the practical implementation of quantum algorithms using modern quantum software frameworks.

\begin{center}
\textbf{How to build, simulate, and execute quantum circuits using industry-standard tools and integrate them with classical machine learning}
\end{center}

\vspace{0.5cm}

\section*{1. QISKIT}

\subsection*{1.1 Overview}

Qiskit is an open-source quantum computing framework developed for building and executing quantum circuits on simulators and real quantum hardware.

\subsection*{1.2 Core Components}

\begin{itemize}
\item \textbf{QuantumCircuit:} circuit construction
\item \textbf{Aer:} simulation backend
\item \textbf{IBM Quantum:} real hardware execution
\end{itemize}

\subsection*{1.3 Circuit Example}

\begin{minted}{python}
from qiskit import QuantumCircuit

qc = QuantumCircuit(1)
qc.h(0)
qc.measure_all()
\end{minted}

\subsection*{1.4 Key Features}

\begin{itemize}
\item Gate-based circuit model
\item Access to real quantum devices
\item Integration with optimization routines
\end{itemize}

\vspace{0.5cm}

\section*{2. PENNYLANE}

\subsection*{2.1 Overview}

PennyLane is a framework designed for quantum machine learning and hybrid quantum-classical computation.

\subsection*{2.2 Core Idea}

Parameterized quantum circuits combined with automatic differentiation:

\[
f(\theta) = \langle \psi(\theta) | H | \psi(\theta) \rangle
\]

\subsection*{2.3 Example}

\begin{minted}{python}
import pennylane as qml

dev = qml.device("default.qubit", wires=1)

@qml.qnode(dev)
def circuit(theta):
    qml.RY(theta, wires=0)
    return qml.expval(qml.PauliZ(0))
\end{minted}

\subsection*{2.4 Key Features}

\begin{itemize}
\item Automatic differentiation
\item Seamless integration with PyTorch, TensorFlow, JAX
\item Native support for variational algorithms
\end{itemize}

\vspace{0.5cm}

\section*{3. CIRQ}

\subsection*{3.1 Overview}

Cirq is a quantum computing framework focused on designing, simulating, and optimizing circuits for near-term quantum devices.

\subsection*{3.2 Circuit Design}

\begin{minted}{python}
import cirq

q = cirq.LineQubit(0)
circuit = cirq.Circuit(
    cirq.H(q),
    cirq.measure(q)
)
\end{minted}

\subsection*{3.3 Key Features}

\begin{itemize}
\item Fine-grained control of quantum circuits
\item Focus on NISQ (Noisy Intermediate-Scale Quantum) devices
\item Hardware-aware circuit optimization
\end{itemize}

\vspace{0.5cm}

\section*{4. APPLICATIONS}

\subsection*{4.1 Circuit Simulation}

Simulation allows testing quantum algorithms without real hardware:

\begin{itemize}
\item Statevector simulation
\item Shot-based simulation
\item Noise modeling
\end{itemize}

Mathematically:

\[
|\psi_{\text{out}}\rangle = U |\psi_{\text{in}}\rangle
\]

\vspace{0.3cm}

\subsection*{4.2 Execution on Quantum Hardware}

\begin{itemize}
\item Access via cloud platforms
\item Limited qubits and noise constraints
\item Measurement statistics from repeated runs
\end{itemize}

Measurement outcome:

\[
P(i) = |\langle i | \psi \rangle|^2
\]

\vspace{0.3cm}

\subsection*{4.3 Integration with Classical Machine Learning}

Hybrid quantum-classical workflows:

\begin{itemize}
\item Quantum circuit as a differentiable layer
\item Classical optimizer updates parameters
\item Feedback loop between quantum and classical systems
\end{itemize}

Optimization loop:

\[
\theta_{k+1} = \theta_k - \eta \nabla_\theta f(\theta)
\]

\vspace{0.5cm}

\section*{FINAL CONNECTION}

All frameworks enable three fundamental capabilities:

\begin{enumerate}
\item \textbf{Construction:} build quantum circuits
\item \textbf{Simulation:} test algorithms classically
\item \textbf{Execution:} run on real quantum hardware
\end{enumerate}

They bridge:
\begin{itemize}
\item Quantum physics
\item Numerical computation
\item Machine learning
\end{itemize}

\vspace{0.5cm}

\section*{STUDY ROADMAP}

\subsection*{Phase 1 --- Qiskit}
\begin{itemize}
\item Basic circuits
\item Simulation
\item Hardware execution
\end{itemize}

\subsection*{Phase 2 --- PennyLane}
\begin{itemize}
\item Parametrized circuits
\item Gradient computation
\item Hybrid models
\end{itemize}

\subsection*{Phase 3 --- Cirq}
\begin{itemize}
\item Circuit optimization
\item Noise modeling
\item Hardware-aware design
\end{itemize}

\subsection*{Phase 4 --- Applications}
\begin{itemize}
\item Quantum simulation
\item Variational algorithms
\item Quantum machine learning
\end{itemize}